\documentclass[11pt]{article}
\usepackage{amsmath,amssymb,amsthm}
\usepackage{algorithm}
\usepackage[noend]{algpseudocode}
\usepackage{graphicx}
\usepackage{hyperref}
\usepackage{bm}
\usepackage{enumitem}
\usepackage{geometry}
\geometry{margin=1in}
\newtheorem{theorem}{Theorem}
\newtheorem{lemma}{Lemma}
\newtheorem{assumption}{Assumption}
\newtheorem{remark}{Remark}
\newcommand{\E}{\mathbb{E}}
\newcommand{\R}{\mathbb{R}}
\newcommand{\norm}[1]{\left\lVert#1\right\rVert}
\newcommand{\ip}[2]{\langle #1,#2\rangle}
\newcommand{\grad}{\nabla}
\begin{document}

\section*{Algorithm Name}
\textbf{Differentially Private Stochastic Gradient Descent with Gaussian Noise (DP‑SGD)}  

The algorithm performs standard SGD updates while adding isotropic Gaussian noise to each (clipped) gradient in order to satisfy $(\varepsilon,\delta)$‑differential privacy.

\section*{Mathematical Formulation}
Let $f:\R^{d}\to\R$ be the empirical risk
\[
f(w)=\frac{1}{n}\sum_{i=1}^{n} \ell(w;z_i),
\]
where $\ell(\cdot;z_i)$ is convex and $L$‑smooth for every data point $z_i$.  
At iteration $t\in\{0,1,\dots,T-1\}$:

\begin{enumerate}[label=\arabic*.]
\item \textbf{Per‑sample gradients} $g_t^{(i)}=\grad\ell(w_t;z_i)$.
\item \textbf{Gradient clipping} 
\[
\tilde g_t^{(i)} = g_t^{(i)}\cdot\min\Bigl\{1,\frac{C}{\norm{g_t^{(i)}}}\Bigr\},
\]
where $C>0$ is the clipping norm.
\item \textbf{Mini‑batch average} (batch size $b$)
\[
\bar g_t = \frac{1}{b}\sum_{i\in\mathcal B_t}\tilde g_t^{(i)} .
\]
\item \textbf{Gaussian perturbation}
\[
\xi_t\sim\mathcal N\bigl(0,\sigma^{2}C^{2}\mathbf I_d\bigr),\qquad 
g_t = \bar g_t + \xi_t .
\]
\item \textbf{Parameter update}
\[
w_{t+1}=w_t-\eta_t g_t ,
\]
with step size $\eta_t>0$.
\item \textbf{Privacy accounting} – the composition of the Gaussian mechanism over $T$ iterations is tracked (e.g., via the moments accountant) to guarantee a total privacy budget $(\varepsilon,\delta)$.
\end{enumerate}

\section*{Pseudocode}
\begin{algorithm}[H]
\caption{DP‑SGD with Gaussian Noise}
\begin{algorithmic}[1]
\State \textbf{Input:} data $\{z_i\}_{i=1}^n$, loss $\ell$, clipping norm $C$, noise scale $\sigma$, learning‑rate schedule $\{\eta_t\}_{t=0}^{T-1}$, batch size $b$, iterations $T$.
\State Initialize $w_0\in\R^{d}$.
\For{$t=0$ \textbf{to} $T-1$}
    \State Sample a minibatch $\mathcal B_t\subseteq\{1,\dots,n\}$, $|\mathcal B_t|=b$.
    \For{each $i\in\mathcal B_t$}
        \State $g_t^{(i)}\gets\grad\ell(w_t;z_i)$
        \State $\tilde g_t^{(i)}\gets g_t^{(i)}\cdot\min\{1,\frac{C}{\norm{g_t^{(i)}}}\}$
    \EndFor
    \State $\bar g_t \gets \frac{1}{b}\sum_{i\in\mathcal B_t}\tilde g_t^{(i)}$
    \State Sample $\xi_t\sim\mathcal N(0,\sigma^{2}C^{2}\mathbf I_d)$
    \State $g_t\gets \bar g_t+\xi_t$
    \State $w_{t+1}\gets w_t-\eta_t g_t$
    \State \textbf{(optional)} Update privacy accountant.
\EndFor
\State \textbf{Output:} $w_T$ (or the averaged iterate $\bar w_T=\frac1T\sum_{t=0}^{T-1}w_t$).
\end{algorithmic}
\end{algorithm}

\section*{Dry Run Example}
Consider the one‑dimensional quadratic $f(w)=(w-3)^2$, thus $\grad f(w)=2(w-3)$.  
Take $C=10$ (no clipping), $\sigma=0$ (to illustrate deterministic dynamics), learning rate $\eta=0.1$, and $w_0=0$.  

\begin{center}
\begin{tabular}{c|c|c|c}
$t$ & $g_t$ (gradient) & $w_{t+1}=w_t-\eta g_t$ & $f(w_{t+1})$\\\hline
0 & $2(0-3)=-6$ & $0-0.1(-6)=0.6$ & $(0.6-3)^2=5.76$\\
1 & $2(0.6-3)=-4.8$ & $0.6-0.1(-4.8)=1.08$ & $(1.08-3)^2=3.6864$\\
2 & $2(1.08-3)=-3.84$ & $1.08-0.1(-3.84)=1.464$ & $(1.464-3)^2=2.3665$\\
\end{tabular}
\end{center}
The objective value decreases monotonically, illustrating the basic behaviour of the update rule. With $\sigma>0$, each $g_t$ would be perturbed by a Gaussian sample, causing stochastic fluctuations while preserving an expected descent.

\newpage
\section*{Assumptions}
\begin{assumption}[Convexity and Smoothness]\label{ass:convex-smooth}
The loss $\ell(\cdot;z)$ is convex and $L$‑smooth for every $z$, i.e.,
\[
\ell(y;z)\ge \ell(x;z)+\ip{\grad\ell(x;z)}{y-x},\qquad
\|\grad\ell(y;z)-\grad\ell(x;z)\|\le L\|y-x\|
\]
for all $x,y\in\R^{d}$.
\end{assumption}
\begin{assumption}[Bounded Gradient Norm after Clipping]\label{ass:bounded}
After clipping, $\norm{\tilde g_t^{(i)}}\le C$ almost surely, thus $\norm{\bar g_t}\le C$.
\end{assumption}
\begin{assumption}[Noise Distribution]\label{ass:noise}
The added noise satisfies $\xi_t\sim\mathcal N(0,\sigma^{2}C^{2}\mathbf I_d)$ and is independent of all past randomness.
\end{assumption}
\begin{assumption}[Learning‑rate Schedule]\label{ass:lr}
The stepsizes are non‑increasing and satisfy
\[
\eta_t\le \frac{1}{L},\qquad \sum_{t=0}^{T-1}\eta_t = \Theta(\sqrt{T}),\qquad
\sum_{t=0}^{T-1}\eta_t^{2}=O(1).
\]
A concrete choice is $\eta_t=\frac{c}{\sqrt{t+1}}$ with $c\le \frac{1}{L}$.
\end{assumption}

\section*{Main Convergence Theorem}
\begin{theorem}[Expected Suboptimality of DP‑SGD]\label{thm:main}
Under Assumptions \ref{ass:convex-smooth}–\ref{ass:lr}, let $w^{*}\in\arg\min_{w}f(w)$. Define the averaged iterate
\[
\bar w_T:=\frac{1}{\sum_{t=0}^{T-1}\eta_t}\sum_{t=0}^{T-1}\eta_t w_t .
\]
Then
\[
\E\!\bigl[f(\bar w_T)-f(w^{*})\bigr]\;\le\;
\frac{\norm{w_0-w^{*}}^{2}}{2\sum_{t=0}^{T-1}\eta_t}
\;+\;
\frac{C^{2}}{2}\,\frac{\sum_{t=0}^{T-1}\eta_t}{\sum_{t=0}^{T-1}\eta_t}
\;+\;
\frac{dC^{2}\sigma^{2}}{2}\,\frac{\sum_{t=0}^{T-1}\eta_t}{\sum_{t=0}^{T-1}\eta_t}
\;=\;
O\!\Bigl(\frac{1}{\sqrt{T}}\Bigr).
\]
In particular, for $\eta_t=c/\sqrt{t+1}$,
\[
\E\!\bigl[f(\bar w_T)-f(w^{*})\bigr]\le
\frac{\norm{w_0-w^{*}}^{2}}{2c}\,\frac{1}{\sqrt{T}}
\;+\;
\frac{cC^{2}}{2}\,\frac{1}{\sqrt{T}}
\;+\;
\frac{c d C^{2}\sigma^{2}}{2}\,\frac{1}{\sqrt{T}} .
\]
\end{theorem}

\section*{Theoretical Properties}
\begin{itemize}
\item \textbf{Stability:} The clipping bound $C$ together with the Gaussian perturbation guarantees that the update is $L$‑Lipschitz in $w_t$, preventing exploding iterates.
\item \textbf{Hyper‑parameter sensitivity:}
  \begin{itemize}
  \item Larger $C$ reduces bias from clipping but inflates the noise variance $dC^{2}\sigma^{2}$.
  \item The stepsize constant $c$ trades off the $1/\sqrt{T}$ coefficient: too large $c$ violates the $1/L$ bound, too small $c$ slows convergence.
  \end{itemize}
\item \textbf{Comparison to non‑private SGD:} Setting $\sigma=0$ and $C$ sufficiently large recovers the classic $O(1/\sqrt{T})$ rate for convex $L$‑smooth functions. The additional term $dC^{2}\sigma^{2}$ quantifies the privacy‑induced error.
\item \textbf{Privacy‑utility trade‑off:} For a fixed privacy budget, $\sigma$ is essentially $\Theta(1/\sqrt{n})$; thus the excess error scales as $O(dC^{2}/(n\sqrt{T}))$.
\end{itemize}

\section*{Discussion}
The theorem shows that DP‑SGD retains the optimal $O(1/\sqrt{T})$ convergence rate for convex smooth objectives, up to an additive variance term stemming from the Gaussian noise required for differential privacy. By carefully calibrating $C$, $\sigma$, and the learning‑rate schedule, the privacy‑induced penalty can be made comparable to the stochastic gradient variance already present in standard SGD.

\newpage
\section*{Preliminaries (Definitions and Lemmas)}
\begin{lemma}[Smoothness Inequality]\label{lem:smooth}
If $f$ is $L$‑smooth, then for any $x,y\in\R^{d}$,
\[
f(y)\le f(x)+\ip{\grad f(x)}{y-x}+\frac{L}{2}\norm{y-x}^{2}.
\]
\end{lemma}
\begin{proof}
Standard from the integral form of the gradient; omitted for brevity. \hfill $\square$
\end{proof}
\begin{lemma}[Unbiasedness of the Noisy Gradient]\label{lem:unbiased}
Conditioned on $w_t$, the noisy gradient satisfies
\[
\E[g_t\mid w_t]=\E[\bar g_t\mid w_t]+\E[\xi_t]=\E[\bar g_t\mid w_t],
\]
and $\E[\xi_t]=0$. Moreover $\E[\norm{\xi_t}^{2}]=dC^{2}\sigma^{2}$.
\end{lemma}
\begin{proof}
Independence of $\xi_t$ from $w_t$ and zero mean of the Gaussian give the result. \hfill $\square$
\end{proof}

\section*{Main Proof (Step‑by‑Step Derivation)}
We prove Theorem \ref{thm:main}. Throughout, denote $w^{*}\in\arg\min f$.

\begin{enumerate}[label={[\textbf{Step \arabic*}]}]
\item \textbf{Expand the squared distance after one update.}
\[
\begin{aligned}
\norm{w_{t+1}-w^{*}}^{2}
&=\norm{w_t-\eta_t g_t-w^{*}}^{2}\\
&=\norm{w_t-w^{*}}^{2}
-2\eta_t\ip{g_t}{w_t-w^{*}}
+\eta_t^{2}\norm{g_t}^{2}.
\end{aligned}
\tag{1}
\]

\item \textbf{Take conditional expectation w.r.t. $w_t$.} Using Lemma \ref{lem:unbiased},
\[
\E\!\bigl[\norm{w_{t+1}-w^{*}}^{2}\mid w_t\bigr]
= \norm{w_t-w^{*}}^{2}
-2\eta_t\ip{\E[g_t\mid w_t]}{w_t-w^{*}}
+\eta_t^{2}\E\!\bigl[\norm{g_t}^{2}\mid w_t\bigr].
\tag{2}
\]

\item \textbf{Bound the inner product by convexity.} Since $f$ is convex,
\[
f(w_t)-f(w^{*})\le \ip{\grad f(w_t)}{w_t-w^{*}}.
\tag{3}
\]
Because $\E[\bar g_t\mid w_t]=\grad f(w_t)$ (the minibatch average is unbiased for the full gradient), we obtain
\[
\ip{\E[g_t\mid w_t]}{w_t-w^{*}}=\ip{\grad f(w_t)}{w_t-w^{*}}
\ge f(w_t)-f(w^{*}).
\tag{4}
\]

\item \textbf{Bound the second moment of $g_t$.}
\[
\begin{aligned}
\E\!\bigl[\norm{g_t}^{2}\mid w_t\bigr]
&= \E\!\bigl[\norm{\bar g_t+\xi_t}^{2}\mid w_t\bigr] \\
&= \E\!\bigl[\norm{\bar g_t}^{2}\mid w_t\bigr]
+2\E\!\bigl[\ip{\bar g_t}{\xi_t}\mid w_t\bigr]
+\E\!\bigl[\norm{\xi_t}^{2}\bigr] .
\end{aligned}
\tag{5}
\]
Independence of $\xi_t$ and zero mean give the cross term zero, and by Assumption \ref{ass:bounded} $\norm{\bar g_t}\le C$, hence
\[
\E\!\bigl[\norm{\bar g_t}^{2}\mid w_t\bigr]\le C^{2},\qquad
\E\!\bigl[\norm{\xi_t}^{2}\bigr]=dC^{2}\sigma^{2}.
\tag{6}
\]
Thus
\[
\E\!\bigl[\norm{g_t}^{2}\mid w_t\bigr]\le C^{2}+dC^{2}\sigma^{2}.
\tag{7}
\]

\item \textbf{Insert bounds (4) and (7) into (2).}
\[
\begin{aligned}
\E\!\bigl[\norm{w_{t+1}-w^{*}}^{2}\mid w_t\bigr]
&\le \norm{w_t-w^{*}}^{2}
-2\eta_t\bigl[f(w_t)-f(w^{*})\bigr]
+\eta_t^{2}\bigl(C^{2}+dC^{2}\sigma^{2}\bigr).
\end{aligned}
\tag{8}
\]

\item \textbf{Rearrange to isolate the suboptimality term.}
\[
2\eta_t\bigl[f(w_t)-f(w^{*})\bigr]
\le \norm{w_t-w^{*}}^{2}
-\E\!\bigl[\norm{w_{t+1}-w^{*}}^{2}\mid w_t\bigr]
+\eta_t^{2}\bigl(C^{2}+dC^{2}\sigma^{2}\bigr).
\tag{9}
\]

\item \textbf{Take full expectation and sum over $t=0,\dots,T-1$.}
Denote $\E_t[\cdot]=\E[\cdot\mid w_t]$ and $\E[\cdot]$ the total expectation.
Summing (9) and using linearity,
\[
\begin{aligned}
2\sum_{t=0}^{T-1}\eta_t\,\E\!\bigl[f(w_t)-f(w^{*})\bigr]
&\le \norm{w_0-w^{*}}^{2}
-\E\!\bigl[\norm{w_T-w^{*}}^{2}\bigr]
+\sum_{t=0}^{T-1}\eta_t^{2}\bigl(C^{2}+dC^{2}\sigma^{2}\bigr) .
\end{aligned}
\tag{10}
\]

\item \textbf{Drop the non‑negative term $\E[\norm{w_T-w^{*}}^{2}]$.}
\[
2\sum_{t=0}^{T-1}\eta_t\,\E\!\bigl[f(w_t)-f(w^{*})\bigr]
\le \norm{w_0-w^{*}}^{2}
+\bigl(C^{2}+dC^{2}\sigma^{2}\bigr)\sum_{t=0}^{T-1}\eta_t^{2}.
\tag{11}
\]

\item \textbf{Define the weighted average iterate.}
Let $S_T:=\sum_{t=0}^{T-1}\eta_t$ and
\[
\bar w_T:=\frac{1}{S_T}\sum_{t=0}^{T-1}\eta_t w_t .
\]
Convexity of $f$ yields
\[
f(\bar w_T)-f(w^{*})\le\frac{1}{S_T}\sum_{t=0}^{T-1}\eta_t\bigl[f(w_t)-f(w^{*})\bigr].
\tag{12}
\]

\item \textbf{Combine (11) and (12).}
\[
\begin{aligned}
\E\!\bigl[f(\bar w_T)-f(w^{*})\bigr]
&\le\frac{1}{2S_T}\Bigl(\norm{w_0-w^{*}}^{2}
+\bigl(C^{2}+dC^{2}\sigma^{2}\bigr)\sum_{t=0}^{T-1}\eta_t^{2}\Bigr).
\end{aligned}
\tag{13}
\]

\item \textbf{Plug in the stepsize schedule $\eta_t=c/\sqrt{t+1}$.}
We have
\[
S_T=\sum_{t=0}^{T-1}\frac{c}{\sqrt{t+1}}\;\ge\;2c\bigl(\sqrt{T}-1\bigr)=\Theta(c\sqrt{T}),
\]
and
\[
\sum_{t=0}^{T-1}\eta_t^{2}=c^{2}\sum_{t=0}^{T-1}\frac{1}{t+1}=c^{2}(\log T+O(1))=O(c^{2}\log T).
\]
Since $\log T = o(\sqrt{T})$, the dominant term in (13) is $\frac{\norm{w_0-w^{*}}^{2}}{2S_T}=O\!\bigl(\frac{1}{\sqrt{T}}\bigr)$. Keeping the exact constants gives
\[
\E\!\bigl[f(\bar w_T)-f(w^{*})\bigr]
\le
\frac{\norm{w_0-w^{*}}^{2}}{2c}\frac{1}{\sqrt{T}}
\;+\;
\frac{c\bigl(C^{2}+dC^{2}\sigma^{2}\bigr)}{2}\frac{1}{\sqrt{T}}
\;+\;o\!\bigl(T^{-1/2}\bigr).
\tag{14}
\]

\item \textbf{Conclusion.} Inequality (14) is exactly the bound stated in Theorem \ref{thm:main}. \hfill $\blacksquare$
\end{enumerate}

\section*{Rate Derivation (With Constants)}
From (13) and the generic bound $\sum_{t=0}^{T-1}\eta_t^{2}\le \eta_{\max}\sum_{t=0}^{T-1}\eta_t$ (valid for non‑increasing $\eta_t$), we obtain a clean $O(1/\sqrt{T})$ expression:
\[
\E\!\bigl[f(\bar w_T)-f(w^{*})\bigr]
\le
\frac{\norm{w_0-w^{*}}^{2}}{2S_T}
+\frac{\eta_{\max}\bigl(C^{2}+dC^{2}\sigma^{2}\bigr)}{2}
\frac{S_T}{S_T^{2}}
=
\frac{\norm{w_0-w^{*}}^{2}}{2S_T}
+\frac{\eta_{\max}\bigl(C^{2}+dC^{2}\sigma^{2}\bigr)}{2S_T}.
\]
With $\eta_{\max}=c$, $S_T\ge 2c\sqrt{T-1}$, we finally have
\[
\E\!\bigl[f(\bar w_T)-f(w^{*})\bigr]\le
\frac{\norm{w_0-w^{*}}^{2}+c^{2}\bigl(C^{2}+dC^{2}\sigma^{2}\bigr)}{4c\sqrt{T-1}}
=O\!\Bigl(\frac{1}{\sqrt{T}}\Bigr).
\]

\section*{Special Cases}
\begin{enumerate}[label=\arabic*.]
\item \textbf{No privacy noise} ($\sigma=0$). The bound reduces to the classical SGD rate with clipping bias $C$.
\item \textbf{High‑dimensional regime} ($d\gg 1$). The excess term scales linearly with $d$, reflecting the fact that isotropic Gaussian noise adds variance in each coordinate.
\item \textbf{Strongly convex $f$}. If $f$ is $\mu$‑strongly convex, a similar derivation (replacing convexity inequality (3) by strong convexity) yields an $O(\frac{\log T}{T})$ rate, with an additive $O(dC^{2}\sigma^{2})$ floor.
\end{enumerate}

\end{document}